\providecommand{\Hilb}{\mathrm{Hilb}}
\providecommand{\Cx}{\mathbb C^{\times}}
\providecommand{\ch}{\mathrm{ch}}
\providecommand{\cat}{\mathrm{cat}}
\providecommand{\BKM}{\mathrm{BKM}}
\providecommand{\kch}{\kappa_{\ch}}
\providecommand{\kcat}{\kappa_{\cat}}
\providecommand{\kBKM}{\kappa_{\BKM}}

\section*{$\Hilb^{n}(K3\times E)$ and residual $\Cx_{E}$-equivariance}

\subsection*{Five attack--heal cycles}

\paragraph{Cycle~1 — smoothness of $\Hilb^{n}(K3\times E)$.}
\textit{Attack.} $K3\times E$ has complex dimension three. Fogarty's
theorem is a surface statement, so smoothness of the punctual Hilbert
scheme on a threefold is \emph{not} automatic; on generic threefolds
$\Hilb^{n}$ is singular from length~$4$ onward (Iarrobino).
\textit{Heal.} The object of interest is not the nested Hilbert scheme
of length-$n$ zero-dimensional subschemes of the threefold, but the
\emph{fibrewise} Hilbert scheme
\[
  \Hilb^{n}(K3\times E) \;=\;
  \Hilb^{n}_{K3/E}\bigl(K3\times E \xrightarrow{\,\mathrm{pr}_{E}\,} E\bigr),
\]
parameterising length-$n$ subschemes supported on fibres of
$\mathrm{pr}_{E}$. Each fibre is a K3 surface, on which Fogarty 1968
gives smoothness of $\Hilb^{n}(K3)$ of dimension $2n$; smoothness is
preserved under the smooth base-change $\mathrm{pr}_{E}\colon K3\times
E\to E$, so
\[
  \Hilb^{n}(K3\times E) \;\longrightarrow\; E
\]
is a smooth proper morphism with fibre $\Hilb^{n}(K3)$, total space
smooth of complex dimension $2n+1$. \emph{(Beauville–Donagi applies
to the original K3 fibre, not to the relative Hilbert scheme of the
threefold; Hilbert-of-threefold singularities are bypassed by the
relative-over-$E$ definition adopted here, which is the one used by
Oberdieck–Pixton.)}

\paragraph{Cycle~2 — the $\Cx_{E}$-action and its fixed locus.}
\textit{Attack.} The elliptic curve $E$ has no $\Cx$-action compatible
with its group structure: $\Cx$ acts on the Lie algebra of $E$ by
scaling but does not integrate to an $E$-action, so there is no literal
$\Cx$-action on $E$ or on $K3\times E$.
\textit{Heal.} Pick an \emph{$E$-translation} generator $t\in E$ of
infinite order and let $\Cx_{E}\subset \mathrm{Aut}(E)$ be the
one-parameter subgroup of the \emph{universal cover} $\mathbb C\to E$
acting on the relative cohomology via the infinitesimal
translation~$\partial_{z}$. Equivariant cohomology is taken with
respect to the torus $T_{E}$ generated by the $E$-translation action on
$\Hilb^{n}(K3\times E)/E$ after passing to the universal-cover-blown-up
model $\mathbb C \times K3 \to E$, in the Maulik–Okounkov sense. The
\emph{fixed locus} under this residual action is
\[
  \Hilb^{n}(K3\times E)^{\Cx_{E}}
   \;=\; \coprod_{m\,\mid\,n} \Hilb^{n/m}(K3)\times E[m],
\]
the torsion-section stratification: fixed subschemes are pulled back
from $\Hilb^{n/m}(K3)$ along the $m$-torsion section $E[m]\hookrightarrow E$
(Oberdieck–Pixton 2017, arXiv:1706.10100, §2.3).

\paragraph{Cycle~3 — equivariant cohomology.}
By localisation (Atiyah–Bott, in the Maulik–Okounkov variant for
symplectic resolutions; Nakajima 1997 Chapter~7 for the K3-fibre case),
\[
  H^{*}_{\Cx_{E}}\bigl(\Hilb^{n}(K3\times E)\bigr)
   \;\simeq\;
   \bigoplus_{m\,\mid\,n}\,
     H^{*}\bigl(\Hilb^{n/m}(K3)\bigr)\otimes
     H^{*}_{\Cx_{E}}\bigl(E[m]\bigr),
\]
with the equivariant parameter $\hbar_{E}\in H^{*}_{\Cx_{E}}(\mathrm{pt})$
inverted in the localised theory. On the fibre $\Hilb^{n/m}(K3)$,
Nakajima's operators $\alpha_{-k}$ furnish the Fock-space generators;
the $E$-torsion factor adjoins the level-$m$ Heisenberg refinement
predicted by Oberdieck–Pixton.

\paragraph{Cycle~4 — Göttsche-type product, adapted.}
\textit{Attack.} The naive adaptation
\[
  \sum_{n\geq 0} q^{n}\,\chi\bigl(\Hilb^{n}(K3\times E)\bigr)
   \stackrel{?}{=}\prod_{m\geq 1}\frac{1}{(1-q^{m})^{24}}
\]
is \emph{wrong}: $\chi(K3\times E)=24\cdot 0=0$, so the threefold Euler
number vanishes, the na\"{i}ve Göttsche exponent drops to zero, and
the generating function collapses to~$1$. The Euler characteristic
detects nothing.
\textit{Heal.} Replace $\chi$ with the $\Cx_{E}$-\emph{equivariant
Euler characteristic} $\chi_{\Cx_{E}}$ (equivalently, the Euler form
of the $E$-equivariant K-theory). Localisation to the $E[m]$-fixed
locus gives
\[
  \sum_{n\geq 0} q^{n}\,\chi_{\Cx_{E}}\bigl(\Hilb^{n}(K3\times E)\bigr)
   \;=\;
  \prod_{m\geq 1}\frac{1}{(1-q^{m})^{24}}
  \;=\; \frac{q}{\eta(q)^{24}},
\]
i.e.\ \emph{Göttsche's K3 product survives} on the $E$-reduced
equivariant theory, pulled back along the torsion-section
stratification of Cycle~2 (Göttsche 1990 \emph{Math.\ Ann.} Theorem~1
on $\Hilb^{n}(K3)$; Oberdieck–Pixton 2017 §2.3 for the $E$-lift by
$m$-torsion sections; Nakajima 1997 Chapter~9 for the Fock-space
structure).

\paragraph{Cycle~5 — match to $\chi(Y^{+}(K3\times E))$.}
\textit{Attack.} The positive-half Yangian
$Y^{+}(K3\times E)$ has been computed in two independent ways; the
numerical comparison is where the identification lives or dies.
\textit{Heal.} On the Oberdieck–Pixton component, the positive-half
CoHA character is
\[
  \chi\bigl(Y^{+}(K3\times E)\bigr)(q,\zeta,p)
   \;=\; \frac{1}{\Phi_{1}(q,\zeta,p)}
   \;=\; \frac{1}{\Delta_{5}(q,\zeta,p)},
\]
with $\Delta_{5}$ the Gritsenko cusp form of weight $5$ on
$\mathrm{Sp}_{4}(\mathbb Z)$. Its $\zeta\to 1,\,p\to 0$ specialisation
recovers $q/\eta(q)^{24}$, matching the Hilbert-scheme generating
function of Cycle~4. The double-specialised identity
\[
  \lim_{\zeta\to 1}\lim_{p\to 0}
   \chi\bigl(Y^{+}(K3\times E)\bigr)(q,\zeta,p)
   \;=\;
   \sum_{n\geq 0} q^{n}\,\chi_{\Cx_{E}}\bigl(\Hilb^{n}(K3\times E)\bigr)
\]
is the chain-level witness of the Oberdieck–Pixton identification on
the K3-fibre--$E$-base specialisation; the full trivariate form
belongs to the Igusa/Gritsenko denominator, not to the Hilbert-scheme
count alone (Oberdieck 2018 \emph{Geom.\ Topol.}~\textbf{22} Theorem~$1$;
Oberdieck 2024 arXiv:2405.03418 Theorem~$1.1$).

\subsection*{$\kappa_{\bullet}$-discipline}

$\kcat(K3\times E)=\chi(\mathcal O_{K3})\cdot\chi(\mathcal O_{E})
=2\cdot 0=0$ on the total space, Künneth-multiplicative; the $E$-base
vanishing drives the threefold Euler number to zero and forces the
Cycle~4 escape through $\Cx_{E}$-equivariance. The Hilbert-scheme
generating function lives at the reduced-equivariant level, not at
the Euler-characteristic level. The Borcherds weight is
$\kBKM(\Phi_{1})=c_{1}(0)/2=5$ via Gritsenko $\Delta_{5}$, and
$\kch(K3\times E)$ is computed on the chiral side through $\Phi$, not
by Hilbert-scheme counting: the four $\kappa_{\bullet}$ values
$\{2,3,5,24\}$ are four \emph{distinct constructions} applied to the
same geometric datum.

\subsection*{Primary literature}

Nakajima 1997 \emph{Lectures on Hilbert Schemes of Points on Surfaces}
Chapters~7, 9; Göttsche 1990 \emph{Math.\ Ann.}~\textbf{286} Theorem~1;
Fogarty 1968 \emph{Amer.\ J.\ Math.}~\textbf{90}; Beauville 1983
\emph{J.\ Diff.\ Geom.}~\textbf{18}; Oberdieck–Pixton 2017
arXiv:1706.10100 §0.4, §2.3; Oberdieck 2018 \emph{Geom.\ Topol.}~\textbf{22}
Theorem~$1$; Oberdieck 2024 arXiv:2405.03418 Theorem~$1.1$;
Gritsenko–Nikulin 1998 \emph{Internat.\ Math.\ Res.\ Notices} (Igusa
cusp form).
